\documentclass{rvcepaper}
\begin{document}

% ---------- HS351TA Entrepreneurship & IPR (CIE-II)
\begin{iempaper}
\paperinfo{shownote=0,date={15\textsuperscript{th} Dec 2025},exam={CIE - II},maxmarks={10 + 50},sem={V},prog={UG},duration={30 + 90 Min},
 title={ENTREPRENEURSHIP \& INTELLECTUAL PROPERTY RIGHTS},code={HS351TA}}
\iemhead
\ptitle{Part A}
\begin{cietable}{M}{BT}{CO}
\q{1.}{Distinguish between Proof of Concept (PoC) and Prototype with respect to risk mitigation.}{2}{1}{3}
\q{2.}{Why is technical feasibility meaningless without market feasibility? Justify.}{2}{1}{3}
\q{3.}{How does reverse engineering weaken trade secret protection?}{2}{2}{3}
\q{4.}{Why is novelty mandatory but aesthetic appeal critical in industrial design protection?}{2}{3}{3}
\q{5.}{Define Copyright.}{2}{2}{3}
\end{cietable}
\ptitle{Part B}
\begin{cietable}{M}{BT}{CO}
\q{1.}{Explain the elements of a business plan. Describe each component with suitable examples.}{10}{2}{2}
\q{2.}{Discuss the complete Entrepreneurial Opportunity Evaluation Process covering:
\begin{bul}
\item Identification of market opportunities and trends
\item Technical feasibility (Prototype and Proof of Concept)
\item Financial feasibility (Cost estimation, revenue projection, and break-even analysis)
\item Role of cross-disciplinary collaboration in technological innovation
\end{bul}}{10}{3}{3}
\q{3.}{Explain the acts which do not constitute Infringement under the Copyright Law.}{10}{2}{3}
\q{4.}{Explain in detail the complete procedure for obtaining Industrial Design protection in India. Further, critically evaluate:
\begin{bul}
\item Grounds and procedure for revocation of registered designs
\item Types of design infringement
\item Civil and criminal remedies available
\end{bul}
Illustrate your answer with at least one real case study.}{10}{3}{2}
\q{5.}{A software company develops a unique algorithm for data security and keeps it confidential. A former employee illegally discloses this information to a competitor.
\begin{num}
\item Identify the type of intellectual property right involved.
\item Explain the legal remedies available under Indian law.
\item Discuss the importance of trade secret protection with reference to this case.
\end{num}}{10}{3}{3}
\end{cietable}
\end{iempaper}

% ---------- IM352IA Operations Management (CIE - II-A)
\begin{iempaper}
\paperinfo{date={15/12/2025},exam={CIE -- II-A},maxmarks={10 + 50},sem={V},prog={UG},duration={30 + 90 Min},
 title={Operations Management},code={IM352IA}}
\iemhead
\begin{cietable}{M}{BT}{CO}
\qpart{Part -- A}
\q{1.}{In time series models, regular fluctuations occurring within a year due to weather or festivals are called \blank[2cm] variations.}{1}{1}{1}
\q{2.}{The difference between the actual value and the forecasted value is known as the \blank[3cm].}{1}{1}{2}
\q{3.}{Qualitative forecasting methods rely heavily on \blank[2cm] judgment.}{1}{1}{1}
\q{4.}{A hotel predicts weekend occupancy using past booking patterns. This uses a \blank[2cm] time series model.}{1}{2}{1}
\q{5.}{A company uses a senior manager's opinion to forecast sales of a new product. This reflects the \blank[2cm] forecasting method.}{1}{1}{1}
\q{6.}{A production line experiencing a delay at one critical machine identifies it as the \blank[2cm] resource.}{1}{2}{4}
\q{7.}{A production manager selects the job with the shortest processing time to reduce average waiting time. This applies the \blank[2cm] rule.}{1}{2}{3}
\q{8.}{A plant with multiple machines generates several two-machine equivalent sequences to approximate the optimal schedule. This procedure follows the \blank[2cm].}{1}{2}{3}
\q{9.}{A manufacturer with two machines creates an optimal job sequence by choosing the smallest processing time among both machines. This applies \blank[2cm] rule.}{1}{2}{3}
\q{10.}{For two jobs processed on multiple machines, the job with smaller processing time on the first machine is scheduled \blank[2cm] in the Johnson-like rule for job shops.}{1}{2}{3}
\qpart{Part -- B}
\q{1}{Explain the importance of measuring forecast error in Operations Management. Discuss the managerial implications of each measure with suitable examples.}{08}{2}{3}
\q{2}{With suitable example, differentiate between forward scheduling and backward scheduling}{08}{2}{4}
\q{3}{A lathe manufacturer had noted the sales invoices from January to August and noticed the sales during each month as 10, 11, 10, 11, 14, 15, 14 \& 15 units respectively. What would be his estimate of sales for September?\newline
(i) If he has followed 4 month weighted moving average (with weights 0.4, 0.3, 0.2 \& 0.1).\newline
(ii) If he has followed SES with $\alpha=0.2$ and his estimated sales for January is 11.}{12}{4}{3}
\q{4}{Consider the following single machine scheduling problem.
\qtabc{|c|c|c|}{\hline Job, $j$ & Processing time, $t_j$ & Due date, $d_j$\\\hline 1 & 12 & 17\\\hline 2 & 10 & 14\\\hline 3 & 10 & 12\\\hline 4 & 9 & 20\\\hline 5 & 14 & 18\\\hline 6 & 17 & 27\\\hline}
Determine the sequence which will minimize the maximum lateness ($L_{max}$). Also, determine the $L_{max}$ with respect to the optimal sequence.}{10}{4}{4}
\q{5}{Use graphical method to minimize the time needed to process the following jobs on the machines shown (i.e. for each machine, find the job which should be scheduled first). Also, calculate the total time elapsed to complete both jobs.
\qtabc{|l|l|c|c|c|c|c|}{\hline Job 1 & Sequence & A & B & C & D & E\\\cline{2-7} & Time (Hrs). & 3 & 4 & 2 & 6 & 2\\\hline Job 2 & Sequence & B & C & A & D & E\\\cline{2-7} & Time (Hrs) & 5 & 4 & 3 & 2 & 6\\\hline}}{12}{4}{4}
\end{cietable}
\end{iempaper}

% ---------- IM353IA Quality Assurance (CIE-II)
\begin{iempaper}
\paperinfo{date={16\textsuperscript{th} Dec 2025},exam={CIE-II},maxmarks={10 + 50},sem={V},prog={UG},duration={30 + 90 Min},
 title={Quality Assurance},code={IM353IA}}
\iemhead
\begin{cietable}{M}{BT}{CO}
\qpart{Part -- A}
\q{1.}{Explain the purpose of a p-chart.}{2}{3}{01}
\q{2.}{What is AOQ?}{2}{2}{01}
\q{3.}{State the difference between single and double sampling plans.}{2}{2}{02}
\q{4.}{What is AOQL?}{2}{3}{02}
\q{5.}{What are multiple sampling plans?}{2}{2}{02}
\multicolumn{5}{|c|}{6.}\\ \hline
\q{7.}{Draw the type - A OC curve for the single sampling plan:\newline \hspace*{0.5cm}$n=50$, $C=1$, $N=200$.}{10}{2}{3}
\q{2}{Differentiate between C chart and U chart}{10}{3}{3}
\q{3}{Lots of Bushing screws from a single supplier were inspected 100 percent with the following results.
\qtabc{|c|c|c|}{\hline Lot No. & Amount in lot & Number rejected\\\hline 01 & 6,000 & 183\\\hline 02 & 6,000 & 131\\\hline 03 & 10,000 & 184\\\hline 04 & 1,783 & 23\\\hline 05 & 3,089 & 171\\\hline 06 & 5,774 & 65\\\hline 07 & 3,000 & 46\\\hline 08 & 2,724 & 85\\\hline}
Draw an appropriate control chart for the data. What can you conclude on the state of control?}{10}{4}{3}
\q{4}{The following table gives the number of missing rivets at the aircraft final inspection:
\qtab{|c|c|c|c|c|c|}{\hline Sl No & Air craft No. & Number of missing rivets & Sl No & Air craft No. & Number of missing rivets\\\hline 1 & 202 & 07 & 6 & 207 & 11\\\hline 2 & 203 & 12 & 7 & 208 & 06\\\hline 3 & 204 & 13 & 8 & 209 & 19\\\hline 4 & 205 & 18 & 9 & 210 & 04\\\hline 5 & 206 & 17 & 10 & 211 & 26\\\hline}
Establish an appropriate control chart for the purpose.}{10}{4}{4}
\q{5}{Explain the meaning of the following terms:
\qtab{|l|l|l|l|}{\hline i. & Producer's risk & ii. & AQL\\\hline iii. & Consumer's risk & iv. & LTPD\\\hline}}{10}{3}{4}
\end{cietable}
\end{iempaper}

% ---------- IM254TA Financial Accounting & Costing (CIE - II)
\begin{iempaper}
\paperinfo{shownote=0,date={16\textsuperscript{th} Dec 2025},exam={CIE - II},maxmarks={10 + 50},sem={V},prog={UG},duration={30 + 90 Min},
 title={FINANCIAL ACCOUNTING \& COSTING},code={IM254TA}}
\iemhead
{\small Note: Answer all the Questions.}\par\vspace{3pt}
\begin{cietable}{M}{BT}{CO}
\qpart{PART - A}
\q{1.1}{What is meant by Overheads in Costing?}{2}{1}{2}
\q{1.2}{Differentiate between Fixed Costs and Variable Costs.}{2}{3}{2}
\q{1.3}{Identify two sectors for which activity-based costing is suitable.}{2}{3}{4}
\q{1.4}{Comment on how cost accumulation happens in Process Costing with an example.}{2}{3}{4}
\q{1.5}{Job costing is for \blank[1.3cm], \blank[1.3cm] jobs, while batch costing is for \textit{a group of} \blank[1.3cm] \textit{products} treated as one batch.}{2}{2}{3}
\qpart{PART - B}
\q{2.a}{Enumerate the objectives of Cost Accounting.}{05}{2}{3}
\q{2.b}{Discuss the various elements of Costing.}{05}{2}{3}
\q{3}{Prepare Trial Balance for the following Balances as on 31\textsuperscript{st} March 2023
\qtab{|l|r|}{\hline \textbf{Balances} & \multicolumn{1}{c|}{\textbf{\rupee}}\\\hline Cost of Goods Sold & 5,20,000\\\hline Opening Stock & 50,000\\\hline Closing Stock & 50,000\\\hline Salary and Wages & 50,000\\\hline Sales & 8,00,000\\\hline Plant \& Machinery & 2,00,000\\\hline Drawing & 50,000\\\hline Investment & 4,30,000\\\hline Creditors & 1,00,000\\\hline Capital & 4,00,000\\\hline}}{10}{3}{2}
\q{4.}{Prepare Profit and Loss account from the following particulars for the year ended March 31, 2023:\newline
Gross Profit \rupee\ 4,69,000\newline Subletting \rupee\ 10,000\newline Commission Received \rupee\ 40,000\newline Rent, tax, and Insurance \rupee\ 10,000\newline Discount Received \rupee\ 22,000\newline Export Duty \rupee\ 15,000\newline Rent Received \rupee\ 23,000\newline Audit Fees \rupee\ 5,000\newline Bank Charges\rupee\ 8,000\newline Packing Charges(Show Room)\rupee\ 12,000\newline Dividend Received \rupee\ 15,000\newline Depreciation \rupee\ 24,000\newline Administration Expenses \rupee\ 35,000\newline Selling Expenses \rupee\ 25,000\newline Printing and Stationery \rupee\ 15,000\newline Distribution Expenses \rupee\ 28,000\newline Office Expenses \rupee\ 32,000\newline Trade Expenses \rupee\ 50,000\newline General Expenses \rupee\ 25,000\newline Carriage outward \rupee\ 15,000\newline Legal Expenses \rupee\ 5,000\newline Bad Debts \rupee\ 10,000}{10}{4}{3}
\q{5}{Write short notes on any three types of costing methods used in organizations.}{10}{2}{3}
\q{6}{In a manufacturing unit, raw material passes through four processes, I, II, III, and IV and the output of each process is the input for the subsequent process. The losses in the four processes are respectively 25\%, 20\%, 20\% and 16 2/3 \% respectively for I, II, III and IV processes of the input. If the end product at the end of the IV process is 40,000 kg, what is the quantity of raw material required to be fed at the beginning of Process I and the cost of the same at Rs. 5 per kg?}{10}{3}{3}
\end{cietable}
\stars{*****}
\end{iempaper}

% ---------- IM355TBB Enterprise Information Systems (CIE - II)
\begin{iempaper}
\paperinfo{date={17/12/2025},exam={CIE -- II},maxmarks={10 + 50},sem={V},prog={UG},duration={30 + 90 Min},
 title={Enterprise Information Systems},code={IM355TBB},note={Answer all questions}}
\iemhead
\begin{cietable}{M}{BT}{CO}
\qpart{Part -- A}
\q{1.}{Define Robotic Process Automation (RPA) and state its primary focus.}{2}{1}{CO1}
\q{2.}{State the fundamental reason why a company must engage in Business Process Reengineering (BPR) before implementing any major automation initiative.}{2}{2}{CO2}
\q{3.}{Define the term Data Mart and state how it differs in \textbf{scope} from a full Data Warehouse.}{2}{1}{CO3}
\q{4.}{Name the two foundational processes that are typically executed by Data Mining tools to discover new knowledge within a BI system.}{2}{2}{CO3}
\q{5.}{A decision-maker uses a BI tool to analyze sales by clicking on a region, then a city, and finally a specific store. Name this OLAP technique and briefly state its purpose.}{2}{2}{CO3}
\qpart{Part -- B}
\q{1}{List and explain various methods of business process automation.}{10}{2}{CO2}
\q{2}{Define Robotic Process Automation (RPA). Explain how it works. List the uses and benefits.}{10}{2}{CO2}
\q{3}{Explain the core difference between Business Intelligence (BI) and Data Warehousing. Describe why BI is considered ``the technology and practice of applying information to make decisions''.}{10}{2}{CO1}
\q{4}{Describe the four progressive types of Data Analytics used in an EIS. For each type, state the central question it is designed to answer.}{10}{3}{CO3}
\q{5}{Apply the concept of Online Analytical Processing (OLAP) by explaining its primary function. Illustrate how a retail manager would use an OLAP technique like Drilling Down to analyze a regional sales report.}{10}{4}{CO3}
\end{cietable}
\end{iempaper}

% ---------- CH375TGD Green and Hydrogen Technology (Global Elective: Group D)
\begin{gepaper}
\paperinfo{program={CHEMICAL ENGINEERING},title={GREEN AND HYDROGEN TECHNOLOGY},code={CH375TGD},exam={CONTINUOUS INTERNAL EVALUATION - 2},ggroup={GLOBAL ELECTIVE: GROUP D},sem={V},month={DEC 2025},maxmarks={60},gdate={16.12.2025},mtime={13:20 - 15:20},duration={120 Min},
 instr={All the questions are compulsory. No additional booklet will be provided.}}
\gehead
\begin{getable}
\gq{1.1}{What are the outcomes of hydrogen risk assessment?}{02}{01}{01}
\gq{1.2}{List the steps involved in the quantitative hydrogen risk assessment.}{02}{01}{04}
\gq{1.3}{Define Codes of Practice.}{02}{01}{03}
\gq{1.4}{What is PAS?}{02}{01}{02}
\gq{1.5}{List the parameters associated with hydrogen for which the codes and standards are defined. Hint: The approval for these parameters vary between countries.}{02}{01}{02}
\gepart{Part B}
\gq{2}{Explain in detail, how hydrogen gas is stored as compressed gas including its advantages and disadvantages.}{10}{02}{04}
\gq{3}{What are the salient features of storing hydrogen as cryogenic liquid? How is it done and what are the associated benefits and hurdles?}{10}{02}{04}
\gq{4}{Is storing hydrogen as a metal hydride a boon or a bane (annoyance)? Justify with substantial details including the working principle, advantages and disadvantages.}{10}{03}{03}
\gq{5}{How are liquid organic hydrogen carriers beneficial in case of storage and transport of hydrogen? Provide a detailed narration.}{10}{02}{03}
\gq{6}{Give an explanatory note on the codes and standards (national and international) associated with the handling and safety of hydrogen gas.}{10}{03}{01}
\end{getable}
\par\vspace{6pt}\noindent{\fontsize{12}{14}\selectfont Course Outcomes}\par\vspace{2pt}
\noindent\begin{tabular}{|C{0.9cm}|p{15.3cm}|}\hline \multicolumn{2}{|l|}{At the end of the course the student will be able to,}\\\hline
1 & Understand the importance of hydrogen and its use as an energy carrier\\\hline
2 & Explain the production, storage and handling of hydrogen\\\hline
3 & Analyze the need for hydrogen as an alternate fuel and the associated challenges\\\hline
4 & Appraise the importance of safety, regulations and codes\\\hline\end{tabular}
\par\vspace{10pt}
\noindent\begin{tabular}{|C{2.2cm}|*{4}{C{1.5cm}|}*{6}{C{1.4cm}|}}\hline Cos/BTL & CO1 & CO2 & CO3 & CO4 & L1 & L2 & L3 & L4 & L5 & L6\\\hline Marks & 12 & 04 & 22 & 22 & 10 & 30 & 20 & . & . & .\\\hline\end{tabular}
\end{gepaper}

\end{document}
